\documentclass[12pt]{article}
\usepackage[utf8]{inputenc}
\usepackage{graphicx}
\usepackage{amsmath}
\usepackage{amssymb}
\usepackage{mathrsfs}
\usepackage{pgfornament}
\usepackage[many]{tcolorbox}
\usepackage[margin = 0.5in]{geometry}
\usepackage{collectbox}
\usepackage{mdframed}
\usepackage{xcolor}
\usepackage{mathtools}
\DeclarePairedDelimiter\ceil{\lceil}{\rceil}
\DeclarePairedDelimiter\floor{\lfloor}{\rfloor}

\newcommand\textbox[1]{%
  \parbox{.333\textwidth}{#1}%
}

\linespread{2}

\makeatletter
\newcommand{\mybox}{%
    \collectbox{%
        \setlength{\fboxsep}{2pt}%
        \fbox{\BOXCONTENT}%
    }%
}
\makeatother

\usepackage{listings}
\usepackage{color}

\definecolor{dkgreen}{rgb}{0,0.6,0}
\definecolor{gray}{rgb}{0.5,0.5,0.5}
\definecolor{mauve}{rgb}{0.58,0,0.82}

\lstset{frame=tb,
  language=R,
  aboveskip=3mm,
  belowskip=3mm,
  showstringspaces=false,
  columns=flexible,
  basicstyle={\small\ttfamily},
  numbers=none,
  numberstyle=\tiny\color{gray},
  keywordstyle=\color{blue},
  commentstyle=\color{dkgreen},
  stringstyle=\color{mauve},
  breaklines=true,
  breakatwhitespace=true,
  tabsize=3
}
\title{MAT 523 HW 6}
\begin{document}

\begin{center}
\begin{large}
\textbf{MAT 523: Functions of a Real Variable I \\
Homework VI \\}
\end{large}
Nolan R. H. Gagnon \\
\end{center}

\noindent \textbf{(1.)}

\vspace{3mm}

\noindent\mybox{\colorbox{lightgray}{\textbf{Statement}}}\hspace{3mm}\hrulefill

\noindent If $E$ has finite outer measure, then for any $\varepsilon>0$, $E$ is the disjoint union of a finite number of subsets, each of which has outer measure at most $\varepsilon.$

\vspace{5mm}

\noindent\mybox{\colorbox{lightgray}{\textbf{Proof}}}\hspace{3mm}\hrulefill



\hfill$\blacksquare$

\pagebreak

\noindent \textbf{(2.)} 

\vspace{3mm}

\noindent\mybox{\colorbox{lightgray}{\textbf{Statement}}}\hspace{3mm}\hrulefill 

\noindent Let $E_1, E_2$ be measurable sets. Then $$m(E_1\cup E_2) + m(E_1\cap E_2) = m(E_1) + m(E_2).$$ 

\vspace{5mm}

\noindent\mybox{\colorbox{lightgray}{\textbf{Proof}}}\hspace{3mm}\hrulefill

By the definition of a measurable set, we have that 
$$
m(E_1) = m(E_1\cap E_2) + m\left(E_1\cap E_2^{C}\right)
$$
and
$$
m(E_2) = m(E_2\cap E_1) + m\left(E_2\cap E_1^{C}\right).
$$
Adding these equations together yields
\begin{align}
m(E_1) + m(E_2) &= m(E_1\cap E_2) + m(E_2\cap E_1) + m\left(E_1\cap E_2^{C}\right) + m\left(E_2\cap E_1^{C}\right).
\end{align}
Now, we can write $E_1\cup E_2$ in the following way:
\begin{align*}
    E_1\cup E_2 &= (E_1\setminus E_2)\cup (E_2\setminus E_1) \cup (E_1\cap E_2) \\
    &=\left(E_1\cap E_2^{C}\right)\cup \left(E_2\cap E_1^{C}\right)\cup (E_1\cap E_2).
\end{align*}
Since this is a disjoint union, it follows that
\begin{align*}
    m(E_1\cup E_2) &= m\left(E_1\cap E_2^{C}\right) + m\left(E_2 \cap E_1^{C}\right) + m(E_1\cap E_2).
\end{align*}
But then equation (1) becomes 
$$
m(E_1) + m(E_2) = m(E_1\cap E_2) + m(E_1\cup E_2),
$$
as required.

\hfill$\blacksquare$

\end{document}
